% !TeX program = pdflatex
\documentclass[11pt,a4paper]{article}
\usepackage[margin=28mm]{geometry}
\usepackage{amsmath,amssymb,amsthm,mathtools}
\usepackage{microtype}
\usepackage{enumitem}
\usepackage{xcolor,xurl,hyperref,bookmark}
\hypersetup{colorlinks=true,linkcolor=blue!40!black,citecolor=blue!40!black,
 urlcolor=blue!40!black,pdftitle={A Reworking of Mochizuki's Hom Theorem over p-Subtame Fields},
 pdfauthor={}}
\setlength{\emergencystretch}{2em}
\setlist[enumerate]{label=\textup{(\roman*)},leftmargin=*,itemsep=2pt}
\numberwithin{equation}{section}
\theoremstyle{plain}
\newtheorem{theorem}{Theorem}[section]
\newtheorem{proposition}[theorem]{Proposition}
\newtheorem{lemma}[theorem]{Lemma}
\newtheorem{corollary}[theorem]{Corollary}
\theoremstyle{definition}
\newtheorem{definition}[theorem]{Definition}
\theoremstyle{remark}
\newtheorem{remark}[theorem]{Remark}
\newcommand{\Zp}{\mathbf Z_p}
\newcommand{\Qp}{\mathbf Q_p}
\newcommand{\Gm}{\mathbf G_m}
\newcommand{\PP}{\mathbf P}
\newcommand{\cO}{\mathcal O}
\newcommand{\et}{\mathrm{\acute et}}
\newcommand{\etale}{\'etale}
\newcommand{\ab}{\mathrm{ab}}
\newcommand{\nr}{\mathrm{nr}}
\newcommand{\tame}{\mathrm{t}}
\DeclareMathOperator{\Hom}{Hom}
\DeclareMathOperator{\Inn}{Inn}
\DeclareMathOperator{\Gal}{Gal}
\DeclareMathOperator{\Spec}{Spec}
\DeclareMathOperator{\Pic}{Pic}
\DeclareMathOperator{\Lie}{Lie}
\DeclareMathOperator{\Aut}{Aut}
\DeclareMathOperator{\ord}{ord}
\DeclareMathOperator{\trdeg}{trdeg}
\renewcommand{\S}{\ensuremath{\mathsection}}
\newcommand{\stacks}[1]{\href{https://stacks.math.columbia.edu/tag/#1}{\texttt{#1}}}
\title{A Reworking of Mochizuki's Hom Theorem\\over \texorpdfstring{$p$}{p}-Subtame Fields}
\author{}
\date{November 13, 2026}
\begin{document}
\maketitle
\begin{abstract}
We rework Mochizuki's proof of the anabelian Hom theorem, following also
Faltings's exposition of its Picard-theoretic arguments.
Let $p$ be a prime, and let $K$ be a subfield of the maximal tamely ramified
extension of a henselian discretely valued field of mixed characteristic
$(0,p)$. For every set of primes $\Sigma$ containing $p$, we prove that
dominant morphisms from a normal geometrically integral $K$-scheme of finite
type to a hyperbolic curve are determined by, and exhaust, the open
homomorphisms of their geometrically pro-$\Sigma$ arithmetic fundamental
groups over the absolute Galois group of $K$, up to geometric inner
automorphisms. The local argument uses Tate's
theorem on $p$-divisible groups and the relative Picard scheme of a semistable
curve. Approximation of Kummer classes then permits the reconstruction of a
morphism from its action on regular differentials.
\end{abstract}

\section{Introduction}
\label{sec:intro}

Grothendieck's anabelian programme asks whether arithmetic fundamental
groups recover varieties and their morphisms \cite{Grothendieck1997}.
For hyperbolic curves, the early rigidity results of Nakamura for
punctured projective lines \cite{Nakamura1990} were followed by
Tamagawa's Isom theorem for affine curves over fields finitely generated
over $\mathbf Q$ \cite[Theorem~(0.3)]{Tamagawa1997} and Mochizuki's
Isom theorem for proper curves over number fields
\cite[Theorem~A]{Mochizuki1996}. Mochizuki subsequently proved the local
pro-$p$ Hom theorem, which recovers dominant morphisms to hyperbolic
curves from open homomorphisms of arithmetic fundamental groups. Its
ground field is sub-$p$-adic, that is, a subfield of a finitely generated
extension of $\Qp$, and its source is smooth
\cite[Theorem~A and Theorem~16.5]{Mochizuki1999}. The theorem remains true
after replacing the geometric fundamental groups by their maximal
pro-$\Sigma$ quotients, for any set of primes $\Sigma$ containing $p$;
this formulation is stated in \cite[Theorem~A and the subsequent remark,
p.~121]{Mochizuki2003}.

The present paper revisits this proof, with particular attention to the
hypothesis on the ground field. The reconstruction through Jacobians,
Kummer classes, and regular differentials is Mochizuki's. Our presentation
also follows Faltings's account of the Picard-torsor and point-approximation
arguments \cite[\S\S5--6]{Faltings1998}. The use of points on successive
covers has its antecedent in Tamagawa's characterization of geometric
sections \cite[Corollary~(2.10)]{Tamagawa1997}; its adaptation to the local
Hom theorem is developed in \cite[\S\S7--10]{Mochizuki1999}.
We carry out the local comparison over complete discrete valuation rings
and their tamely ramified extensions. This leads to the following class
of fields.

\begin{definition}
\label{def:subtame}
Fix a prime $p$. A field $K$ of characteristic zero is \emph{$p$-subtame} if
there exist a henselian discretely valued field $F$ of mixed characteristic
$(0,p)$ and an embedding
\[
                         K\hookrightarrow F^{\tame},
\]
where $F^{\tame}$ denotes its maximal tamely ramified algebraic extension.
The extension $F^{\tame}/F$ is algebraic; its completion is not used in this
definition. The residue field of $F$ is arbitrary.
\end{definition}

One obtains the same class of fields by requiring $F$ to be complete.
Moreover, the class is closed under subfields and finitely generated
extensions (Lemma~\ref{lem:fields}). In particular it contains the
sub-$p$-adic fields and the generalized sub-$p$-adic fields of
\cite[Definition~4.11]{Mochizuki2003}. Mochizuki's theorem over the latter
class is an Isom theorem \cite[Theorem~4.12]{Mochizuki2003}.
We prove the Hom statement over $p$-subtame fields by the local argument
developed below.

Fix an algebraic closure $\bar K$ of $K$, and put $G_K=\Gal(\bar K/K)$.
For a nonempty set of primes $\Sigma$ and a geometrically connected
$K$-scheme $Z$, write
\[
 \Delta_Z^\Sigma=\pi_1(Z_{\bar K})^{(\Sigma)},\qquad
 \Pi_Z^\Sigma=\pi_1(Z)\big/
       \ker\bigl(\pi_1(Z_{\bar K})\longrightarrow\Delta_Z^\Sigma\bigr).
\]
Here $(\Sigma)$ means the maximal pro-$\Sigma$ quotient. Thus
\begin{equation}
 1\longrightarrow\Delta_Z^\Sigma\longrightarrow\Pi_Z^\Sigma
       \longrightarrow G_K\longrightarrow1.
 \label{eq:arithmetic-sequence}
\end{equation}
All homomorphisms of profinite groups are understood to be continuous.
An open homomorphism is one with open image. We suppress geometric base
points; the quotient by geometric inner automorphisms in the theorem
accounts for their choice. The subscript $G_K$ on $\Hom$ means that the
homomorphisms commute with the projections to $G_K$.
A hyperbolic curve is a smooth separated
geometrically connected curve whose smooth proper compactification has
genus $g$ and whose boundary has $r$ geometric points, with $2g-2+r>0$.

\begin{theorem}[Geometrically pro-$\Sigma$ Hom theorem]
\label{thm:main}
Let $K$ be $p$-subtame, and let $\Sigma$ be any set of primes containing $p$.
Let $X$ be a normal geometrically integral finite-type $K$-scheme and $Y$ a
hyperbolic curve over $K$. Then the natural map
\begin{equation}
 \Hom_K^{\mathrm{dom}}(X,Y)\ \xrightarrow{\ \sim\ }\
 \Hom_{G_K}^{\mathrm{open}}(\Pi_X^\Sigma,\Pi_Y^\Sigma)
                  /\Inn(\Delta_Y^\Sigma)
 \label{eq:main}
\end{equation}
is bijective. The scheme $X$ need not be separated.
\end{theorem}

For $\Sigma=\{p\}$ this is the geometrically pro-$p$ Hom theorem. Taking
$\Sigma$ to be the set of all primes gives the corresponding statement for
full \etale fundamental groups. In both cases the arithmetic quotient in
\eqref{eq:arithmetic-sequence} is the full group $G_K$.

Other extensions of anabelian Isom results use different hypotheses on
the ground field. Hoshi proves an Isom theorem for affine hyperbolic
curves over Kummer-faithful fields under point-theoreticity and
Galois-preservation conditions \cite[Theorem~A]{Hoshi2017}.
Tsujimura proves a relative birational Isom theorem for curves over
arbitrary cyclotomic extensions of number fields
\cite[Corollary~B]{Tsujimura2023}. A relative birational Hom analogue
over generalized sub-$p$-adic fields is posed as a question in
\cite[Introduction, Remark~(a)]{Tsujimura2026}.
Theorem~\ref{thm:main} concerns arithmetic fundamental groups of
finite-type schemes, with their projections to the fixed group $G_K$.
Mochizuki's formulation for pro-varieties also includes function fields
\cite[Definition~16.4 and Theorem~16.5]{Mochizuki1999}.

The change in the ground field requires a local argument valid without
any perfectness assumption on the residue field. We prove the needed
comparison using connected $p$-divisible groups, finite Kummer torsors,
and an integral estimate on differentials. The reconstruction is then
expressed in terms of approximation at each finite level. We also give
an extension criterion for normal sources, including singular ones.
The finite-cover method and the passage from pro-$p$ to pro-$\Sigma$
remain those of the preceding works; the latter is already part of
Mochizuki's formulation cited above.

\subsection*{Outline of the proof}

Consider first a surjective homomorphism between the geometrically pro-$p$
arithmetic fundamental groups of two smooth proper curves, with hyperbolic
target. Abelianization gives a homomorphism between the Tate modules of
their Jacobians. To recover a morphism of curves, one needs first to turn
this homomorphism into a map on regular differentials, and then to show
that the map on differentials respects the equations of the target curve.
This is the order of reconstruction in Mochizuki's proof
\cite[\S\S7--14]{Mochizuki1999}.

The first question is local. After embedding $K$ in a tame extension of a
complete discretely valued field, the Tate-module homomorphism descends to
a finite extension of that field. The reason is that a prime-to-$p$ group
acting on a $\Zp$-lattice is detected by reduction modulo $p$
(Lemma~\ref{lem:tate-descent}). On a regular semistable model, inertia
characterizes the Tate module of the connected formal subgroup of the
Jacobian. Tate's full faithfulness theorem therefore gives an integral
homomorphism of formal groups, whose differential is the required map
\cite[\S4.2, Theorem~4 and Corollary~1]{Tate1967}.
Section~\ref{sec:local} also proves that this homomorphism preserves finite
Kummer classes. The point here is to subtract a point defined over the
algebraic constant field before passing to the formal group. This
subtraction gives both the assertion about Kummer classes and an integral
bound on the error in the differential comparison.

One then applies the Kummer assertion on finite covers of the target.
Over a suitable field with pro-$p$ absolute Galois group, the resulting
sections can be approximated, modulo any prescribed power of $p$, by
sections arising from points (Section~\ref{sec:approximation}). Choose a
generic point of the source in a Gauss-valued extension and a derivation
which preserves its valuation ring. The differential values at the
approximating points converge to the values prescribed by the
Tate-module homomorphism. Ratios of canonical differentials, together
with their first derivatives, generate an affine chart of the target;
the derivatives supply the functions missing from the canonical map
in the hyperelliptic case. Passing to the limit in their polynomial
relations produces a morphism of curves
(Proposition~\ref{prop:reconstruction}). Thus only approximation at each
finite level is needed.

This construction initially takes place over an algebraically closed
extension of $K$. Applying it to finite covers identifies the geometric
outer homomorphism induced by the constructed morphism. Faithfulness of the
action on differentials then gives descent to $K$, and the triviality of
centralizers of open subgroups in the geometric target group identifies the arithmetic
homomorphism as well. These arguments prove the proper-target theorem
for curves in Section~\ref{sec:proper-hom}. For a normal source,
Section~\ref{sec:global} uses a generic curve over a purely transcendental
extension of $K$ which preserves the connectedness of every finite
geometric cover. The rational map so
obtained extends to the source: a positive-dimensional fibre of its
graph would contradict the existence of abelian covers of the target
of arbitrarily large degree. Ramified covers at the cusps exclude the
boundary of an open target. Finally, Section~\ref{sec:sigma} applies the
pro-$p$ theorem on every finite pro-$\Sigma$ cover, as in
\cite[Remark after Theorem~16.5]{Mochizuki1999}, to recover the original
pro-$\Sigma$ homomorphism.

\subsection*{Notation and preliminaries}

Until Section~\ref{sec:sigma}, put $\Pi_Z=\Pi_Z^{\{p\}}$ and
$\Delta_Z=\Delta_Z^{\{p\}}$. For $S$ an abelian variety or $\Gm$ over $K$, put
$T_pS=\varprojlim S[p^n](\bar K)$, with transition map $[p]$, and write
$\kappa_{S,n}:S(L)\to H^1(L,S[p^n])$ for the Kummer map over a field
extension $L/K$.
In particular, $T_pS/p^nT_pS=S[p^n](\bar K)$.

We use the correspondence between finite \etale covers and finite
fundamental-group sets \cite[Expos\'e~V]{SGA1}, the arithmetic exact
sequence \cite[Expos\'e~IX, Th\'eor\`eme~6.1]{SGA1}, and invariance of
geometric fundamental groups under algebraically closed extension in
characteristic zero \cite[Expos\'e~XIII, Corollaire~3.2 and
Remarque~3.1.3]{SGA1}. For a smooth proper curve
$C$ with Jacobian $J_C$, the Abel--Jacobi identifications, formed using a
geometric base point, are
\begin{equation}
 \Delta_C^{\ab}=T_pJ_C,\qquad
 \Lie(J_C)^\vee=H^0(C,\omega_C).
 \label{eq:AJ}
\end{equation}
Both identifications are independent of the base point and compatible with
curve maps, norms, and field extension; see \cite[Propositions~2.2 and~9.1]{Milne1986}.
We use semistable reduction and the standard relative Picard construction,
with references supplied in Section~\ref{sec:local}.

\section{The ground field and descent of Tate homomorphisms}
\label{sec:fields}

The valued overfield in Definition~\ref{def:subtame} serves only for the
local construction. We first choose it with separably closed residue field,
so that its tame Galois group is prime to $p$. We then show that a
homomorphism of Tate modules is defined over a finite valued extension,
even when $K$ itself is not.

\begin{lemma}
\label{lem:fields}
\begin{enumerate}
\item In Definition~\ref{def:subtame}, one may require a complete discretely
valued field $M$ with separably closed residue field. With a uniformizer $\pi$,
one may then write
\[
 M^{\tame}=M(\pi^{1/n}:(n,p)=1).
\]
Every finite quotient of $\Gal(M^{\tame}/M)$ has order prime to $p$.
\item Subfields and finite extensions of $p$-subtame fields are $p$-subtame.
\item Finitely generated extensions of $p$-subtame fields are $p$-subtame.
\end{enumerate}
\end{lemma}
\begin{proof}
For a henselian discretely valued $F$, tame ramification gives
$F^{\tame}=F^{\nr}(\pi^{1/n}:(n,p)=1)$
\cite[Tags~\stacks{04GK}, \stacks{0EXW}]{Stacks}. The completion of
$F^{\nr}$ is a complete discretely valued field $M$ with the same uniformizer
and separably closed residue field. Compatible choices of roots embed
$F^{\tame}$ into $M^{\tame}$. All prime-to-$p$ roots of unity lie in $M$ by
Hensel's lemma, so the finite root-adjoining stages are cyclic of
prime-to-$p$ order. Conversely, a complete discretely valued field is
henselian.

Fix $K\subset M^{\tame}$ and an algebraic closure $\Omega$ of $M$. A finite
extension $K(\alpha)/K$ embeds into $\Omega$, and $M(\alpha)/M$ is finite.
Tame ramification is preserved by base change
\cite[Tag~\stacks{0EXY}]{Stacks}; hence
$M(\alpha)M^{\tame}\subset M(\alpha)^{\tame}$.
This proves (ii).

For independent variables $\mathbf t$, the Gauss completion of
$M(\mathbf t)$ is a complete discretely valued field of mixed
characteristic $(0,p)$. By the same base-change assertion,
$M^{\tame}(\mathbf t)$ embeds into its maximal tame extension.
A finitely generated extension is finite over a purely transcendental one,
proving (iii).
\end{proof}

\begin{lemma}[Finite descent of Tate homomorphisms]
\label{lem:tate-descent}
Fix $K\subset M^{\tame}\subset\Omega$ as above. Let $S_1,S_2$ be
Jacobians or $\Gm$ over $K$, and choose a finite extension $M_0/M$ inside
$M^{\tame}$
over which their defining equations descend. Every $G_K$-equivariant map
$u:T_pS_1\to T_pS_2$ is equivariant over the finite extension
\[
                         M_0(S_1[p],S_2[p]).
\]
The same assertion holds for objects over finite extensions of $K$, using
the same valued algebraic closure $\Omega$.
\end{lemma}
\begin{proof}
The lattice $\Hom_{G_{M^{\tame}}}(T_pS_1,T_pS_2)$ is saturated in
$\Hom_{\Zp}(T_pS_1,T_pS_2)$: if $ph$ is invariant, then
$p(\sigma h-h)=0$ implies that $h$ is invariant. Reduction modulo $p$
therefore embeds its reduction into
$\Hom_{\mathbf F_p}(S_1[p],S_2[p])$. Its Galois action factors through
$\Gal(M^{\tame}/M_0)$, a profinite group all of whose finite quotients have
order prime to $p$.

Such an action on a $\Zp$-lattice is detected modulo $p$. Indeed, if the
lattice has rank $r$, then for every $n$ the kernel of
$\operatorname{GL}_r(\Zp/p^n)\to\operatorname{GL}_r(\mathbf F_p)$ is a
finite $p$-group, so its intersection with the finite prime-to-$p$ image
is trivial. The subgroup fixing both $S_1[p]$ and $S_2[p]$ consequently
acts trivially on every $G_{M^{\tame}}$-equivariant homomorphism.
This applies to $u$, since $K\subset M^{\tame}$.

For a finite extension $K(\alpha)/K$ inside $\Omega$, replace $M$ by
$M(\alpha)$ and use
$K(\alpha)\subset M(\alpha)^{\tame}$.
The residue field remains separably closed, and the same argument applies
inside $\Omega$.
\end{proof}

\section{Finite covers and geometric conjugacy}
\label{sec:covers}

The construction from differentials will first yield a morphism of curves
over an algebraically closed field. To recover the prescribed homomorphism
of fundamental groups, we shall apply that construction on finite covers,
following the comparison preceding \cite[Theorem~14.1]{Mochizuki1999}.
The following profinite facts explain how equality on finite torsors gives
equality up to a single geometric conjugation. They apply to every set of
primes $\Sigma$.

\begin{lemma}
\label{lem:cover-tools}
\begin{enumerate}
\item If $D$ is a closed subgroup of a profinite group $P$, two continuous
homomorphisms $\alpha,\beta:H\to P$ are $D$-conjugate if and only if their
images in every finite quotient of $P$ are conjugate by the image of $D$.
It suffices to test a cofinal system of finite quotients.
\item An open subgroup of $\Pi_Z^\Sigma$ is the relative pro-$\Sigma$
fundamental group of the corresponding connected finite \etale cover,
viewed over its field of constants.
\item For a homomorphism $\theta:P_1\to P_2$ over $G$, where
$1\to D_i\to P_i\to G\to1$ is exact, its image is open if and only if
$\theta(D_1)$ is open in $D_2$.
\end{enumerate}
\end{lemma}
\begin{proof}
For (i), the sets
$\{d\in D:\beta\equiv d\alpha d^{-1}\pmod N\}$, for open normal
$N\triangleleft P$, are nonempty closed subsets of the compact group $D$.
They have the finite intersection property, and their intersection gives
one conjugating element.

For (ii), let $\mathcal D$ be a profinite group and let $\mathcal H$ be
the inverse image of an open subgroup of $\mathcal D^{(\Sigma)}$.
Given an open normal $U\triangleleft\mathcal H$ with finite
$\Sigma$-group quotient, take cores in $\mathcal D$. The quotient of the
core of $\mathcal H$ by the core of $U$ embeds into a finite product of
conjugates of $\mathcal H/U$. Both this quotient and the quotient by the
core of $\mathcal H$ are finite $\Sigma$-groups. Thus
$\ker(\mathcal D\to\mathcal D^{(\Sigma)})\subset U$.
It follows that the maximal pro-$\Sigma$ quotient of $\mathcal H$ is
its image in $\mathcal D^{(\Sigma)}$.
Apply this to the geometric fundamental group and then
use the arithmetic exact sequence.

For (iii), the equalities
$\theta(P_1)\cap D_2=\theta(D_1)$ and $\theta(P_1)D_2=P_2$ give
$[P_2:\theta(P_1)]=[D_2:\theta(D_1)]$.
\end{proof}

In particular, passage to the cover corresponding to the image reduces an
open homomorphism to a surjection.

A finite geometric cover need not be defined over the original ground
field. The next observation allows a prescribed geometric Galois group to
be realized after a finite extension of that field.

\begin{lemma}[Realizing a geometric quotient]
\label{lem:quotient-descent}
Let $1\to D\to P\to\Gamma\to1$ be exact, and let $N\triangleleft D$
be open. After restricting to an open normal subgroup of $\Gamma$, the
quotient map $D\to D/N$ extends to a quotient of the restricted extension.
\end{lemma}
\begin{proof}
Choose an open normal $B\triangleleft P$ with $B\cap D\subset N$.
The subgroup $DB$ is the inverse image of an open normal subgroup of
$\Gamma$. Since $[B,D]\subset B\cap D\subset N$, the subgroup $NB$ is
normal in $DB$, and the inclusion of $D$ induces
$D/N\simeq DB/NB$.
\end{proof}

\begin{lemma}[Arithmetic torsor comparison]
\label{lem:torsor-comparison}
Let $\theta:\Pi_X^\Sigma\to\Pi_Y^\Sigma$ be over $G_K$, and let
$f:X\to Y$ be a $K$-morphism. For an open normal $N\triangleleft\Pi_Y^\Sigma$,
let $V_N\to Y$ be the corresponding $\Pi_Y^\Sigma/N$-torsor, and let
$W_N\to X$ be the pullback by $\theta$. If, for a cofinal system of $N$,
there is a $\Pi_Y^\Sigma/N$-equivariant isomorphism
\[
 W_N\simeq X\times_YV_N
\]
inducing the identity on the common constant-field quotient, then $f_*$ and
$\theta$ differ by one inner automorphism from $\Delta_Y^\Sigma$.
\end{lemma}
\begin{proof}
The isomorphism of finite fibres conjugates the two monodromy actions by an
element of $\Pi_Y^\Sigma/N$. Its identity action on the constant quotient says
that this element lies in the image of $\Delta_Y^\Sigma$. Apply
Lemma~\ref{lem:cover-tools}(i). 
\end{proof}

\section{Integral realization and finite Kummer classes}
\label{sec:local}
\label{sec:local-package}

The passage from Tate modules to formal groups and geometric cohomology
classes is already present in Mochizuki's proof
\cite[\S6 and the proof of Proposition~7.4]{Mochizuki1999}; see also
Faltings's account in \cite[Proposition~11]{Faltings1998}.  We need an
integral version at each finite Kummer level over a discrete valuation
ring with arbitrary residue field.  We obtain it by characterizing the
connected Tate lattice by inertia and applying Tate's full faithfulness
theorem.  The differential of the resulting formal homomorphism will
then supply the map on regular differentials.

Throughout, $R_0$ is a complete discrete valuation ring of mixed
characteristic $(0,p)$, $K_0$ is its fraction field, and $\Omega$ is an
algebraic closure of $K_0$, endowed with the extended valuation.  The
residue field of $R_0$ is arbitrary.

\subsection{Semistable models and their connected torsion}

Let $C/K_0$ be a smooth proper geometrically connected curve.  By
semistable reduction \cite[\S9.2, Theorem~7]{BLR1990}, after a finite
extension of $K_0$ we may choose a regular split semistable model
$\mathcal C/R_0$.  The stable reduction theorem in genus at least two is
due to Deligne and Mumford \cite[Corollary~2.7]{DeligneMumford1969}.
A further finite extension gives a section through the smooth locus
\cite[\S2.3, Proposition~5]{BLR1990}.  Put
\[
 J=\operatorname{Pic}^0_{C/K_0},\qquad
 \mathcal G=\operatorname{Pic}^0_{\mathcal C/R_0}.
\]
Here $\operatorname{Pic}^0$ denotes the subgroup of classes of
multidegree zero.  Raynaud's theorem identifies $\mathcal G$ with the
identity component of the N\'eron model of $J$
\cite[Th\'eor\`emes~8.1.4(b) and~8.2.1]{Raynaud1970}; we use the
formulation in \cite[\S9.5, Theorem~4(b) and Remark~5]{BLR1990}.
It is smooth,
separated, and of finite type, with generic fibre $J$.  Its special fibre
is an extension of the product of the Jacobians of the normalized
components by a split torus
\cite[\S9.2, Example~8]{BLR1990}.
Thus $\mathcal G$ is semiabelian, and its formation commutes with
faithfully flat extensions of discrete valuation rings
\cite[\S7.4, Corollary~4]{BLR1990}.

Write $\widehat{\mathcal G}$ for the formal completion at the identity.
Multiplication by $p^n$ on $\mathcal G$ is faithfully flat and
quasi-finite.  Its local factor at the identity of the special fibre is a
finite flat open and closed subgroup
\[
 \widehat{\mathcal G}[p^n]\subset\mathcal G[p^n].
\]
These kernels form the connected $p$-divisible group associated to
$\widehat{\mathcal G}$
\cite[\S2.2, Proposition~1]{Tate1967}.  We shall also use the same
notation for $S=\mathbf G_m$, with $\mathcal G=\mathbf G_{m,R_0}$,
whose formal $p^n$-torsion is $\mu_{p^n}$.

The following description of $T_p\widehat{\mathcal G}$ makes it
functorial for Galois equivariant homomorphisms of Tate modules.

\begin{lemma}[The connected Tate module]
\label{loc:inertia-lattice}
Let $S$ be one of these Jacobians or $\mathbf G_m$, with model
$\mathcal G$.  Then $T_p\widehat{\mathcal G}$ is saturated in $T_pS$,
the quotient $T_pS/T_p\widehat{\mathcal G}$ is unramified, and
\[
 T_p\widehat{\mathcal G}
 =T_pS\cap\sum_{\sigma\in I_{K_0}}
       (\sigma-1)(T_pS\otimes_{\mathbf Z_p}\mathbf Q_p).
\]
In particular, every $G_{K_0}$-equivariant map between such Tate modules
preserves their connected lattices.  A rationally surjective map is
rationally surjective on the connected lattices as well.
\end{lemma}
\begin{proof}
The injections
$\widehat{\mathcal G}[p^n](\Omega)\hookrightarrow S[p^n](\Omega)$
show that $T_p\widehat{\mathcal G}$ is saturated in $T_pS$.  To prove
that the quotient is unramified, fix $\tau\in J[p^n](\Omega)$ and
$\sigma\in I_{K_0}$.  Represent $\tau$ by a degree-zero divisor, and
choose a finite Galois extension over which all its points are rational.
A point $z$ and its conjugate $\sigma z$ have the same specialization.
If this specialization is smooth, the difference of their horizontal
divisors restricts to the trivial line bundle on the special fibre.
In nodal coordinates, the function
\[
                       \frac{x-x(\sigma z)}{x-x(z)}
\]
defines the same horizontal difference near the node, since the
vertical orders cancel.  On the normalized special branches its
reductions are nonzero constants.  Summing these differences gives an
extension of $\sigma\tau-\tau$ whose restriction to each normalized
component is trivial.  Its $p^n$-th power is trivial by separatedness of
$\mathcal G$.  Its reduction is therefore a $p^n$-torsion point of the
gluing torus, hence the identity in characteristic $p$.  It follows that
$\sigma\tau-\tau\in\widehat{\mathcal G}[p^n](\Omega)$, as required.
For $\mathbf G_m$ the formal and full Tate modules coincide.

It remains to show that the inertia differences span
$T_p\widehat{\mathcal G}\otimes\mathbf Q_p$.
More generally, the rational Tate module of a connected $p$-divisible
group over $R_0$ has no nonzero unramified quotient.  The image of the
integral Tate module in such a quotient is a Galois stable lattice,
hence the Tate module of an \etale $p$-divisible group over $R_0$
\cite[\S1.4]{Tate1967}.  Tate's full faithfulness theorem extends the
quotient map to a homomorphism over $R_0$
\cite[\S4.2, Theorem~4]{Tate1967}.  This homomorphism is zero.  Indeed,
after strict henselization each finite level of the source is connected,
whereas the target is a disjoint union of copies of the base.  The
quotient of $T_p\widehat{\mathcal G}\otimes\mathbf Q_p$ by the span of
its inertia differences is unramified, and must therefore vanish.
Since $T_pS/T_p\widehat{\mathcal G}$ is
unramified, every inertia difference in $T_pS\otimes\mathbf Q_p$ lies
in $T_p\widehat{\mathcal G}\otimes\mathbf Q_p$.  Saturation gives the
displayed formula.  An equivariant homomorphism
preserves inertia differences, and a surjective homomorphism maps their
span onto the corresponding span in the target.  This proves the final
assertions.
\end{proof}

\begin{proposition}[Tate realization]
\label{loc:tate-realization}
Let $S_i$ and $\mathcal G_i$ be as above.  A $G_{K_0}$-equivariant
homomorphism $u:T_pS_1\to T_pS_2$ induces a unique formal group
homomorphism over $R_0$
\[
 F_u:\widehat{\mathcal G}_1\longrightarrow\widehat{\mathcal G}_2,
 \qquad T_pF_u=u|_{T_p\widehat{\mathcal G}_1}.
\]
The assignment $u\mapsto dF_u$ is additive and compatible with
composition and finite extensions.  After base change to $\Omega$, the
induced map on Lie algebras is independent of the semistable models and
finite local stages used in its construction.  For an algebraic
homomorphism $f$ defined after a finite extension of $K_0$, one has
$dF_{T_pf}=df$ over that extension.  Moreover, $dF_u\ne0$ if either
$u\otimes\mathbf Q_p$ is surjective and $\dim S_2>0$, or
$S_1=\mathbf G_m$ and $u\ne0$.
\end{proposition}
\begin{proof}
By Lemma~\ref{loc:inertia-lattice}, $u$ restricts to the connected
lattices.  Tate's theorem extends this restriction uniquely to a
homomorphism of $p$-divisible groups, and hence of their associated
formal groups \cite[\S4.2, Theorem~4 and Corollary~1; \S2.2,
Proposition~1]{Tate1967}.  Uniqueness gives additivity and compatibility
with composition.  The identification with the semistable N\'eron
identity component, and its compatibility with base change, show that
the construction is independent of the chosen model and finite stage.
An algebraic homomorphism extends by the N\'eron mapping property; its
formal completion induces the prescribed Tate map and therefore
coincides with $F_{T_pf}$.  Taking derivatives proves
$dF_{T_pf}=df$.

If $u\otimes\mathbf Q_p$ is surjective, Lemma~\ref{loc:inertia-lattice}
shows that its restriction to the connected part is surjective over
$\mathbf Q_p$.  This target is nonzero when $\dim S_2>0$.  If
$S_1=\mathbf G_m$, the connected Tate module is the whole Tate module.
Thus either hypothesis implies $F_u\ne0$.  The formal logarithm in
characteristic zero then gives $dF_u\ne0$.
\end{proof}

\subsection{Correcting points by algebraic constants}

Mochizuki translates points by constants to place them in the identity
component of a N\'eron model in the argument preceding
\cite[Lemma~9.2]{Mochizuki1999}.  For the valued extensions used here,
we make this correction directly on divisors of the semistable curve.
The correcting point lies in $S(\Omega)$ and is therefore divisible by
every power of $p$; its subtraction does not change a finite Kummer class.

\begin{lemma}[Points modulo constants]
\label{loc:point-correction}
Let $F/K_0$ be an extension of complete discretely valued fields whose
valuation ring dominates $R_0$, and let $E/F$ be algebraic with
$\Omega\subset E$.  For a model $\mathcal G$ as above, every
$P\in S(E)$ admits an expression
\[
                         P=z+c,
 \qquad z\in\mathcal G(\mathcal O_E),\quad c\in S(\Omega).
\]
For $\mathbf G_m$ the displayed sum is interpreted multiplicatively.
\end{lemma}
\begin{proof}
For $S=\mathbf G_m$, choose $c\in\Omega^\times$ with $v(c)=v(P)$;
this is possible since the values are rational multiples of $v(p)$.
Then $P/c\in\mathcal O_E^\times$.

Now suppose $S=J$.  Choose a finite extension $F'/F$ inside $E$ over
which $P$ is defined.  The section of $\mathcal C$ allows us to represent
$P$ by a degree-zero line bundle, and hence by a divisor $\sum m_ZZ$
on $C_{F'}$ \cite[\S8.1, Proposition~4]{BLR1990}.
It suffices to choose constant points $c_Z\in C(\Omega)$ for which
\[
       \overline Z-[F'(Z):F']\,\overline{c_Z}
\]
is Cartier of multidegree zero.  We may enlarge $F'$ by finite constant
extensions inside $E$ to make the chosen $c_Z$ rational.  If such an
extension splits $Z$, the same choice of constant point is retained for
each point above it.

If $Z$ specializes to the smooth locus, take $c_Z$ specializing to a
smooth point of the same component; Hensel's lemma gives such a point
over a finite constant extension.  The closure of $Z$ is Cartier and
has special degree $[F'(Z):F']$ on that component and degree zero on
the others.  The required difference therefore has multidegree zero.

Suppose instead that $Z$ specializes to a node.  In a split \etale
chart the completed local ring is $\mathcal O_{F'}[[x,y]]/(xy-q)$, with
$q\in R_0$.  Normalize $v$ to have value group $\mathbf Z$ on $F'$,
and put $a=x(Z)$.  Then $0<v(a)<v(q)$ and $F'(Z)=F'(a)$: the other
coordinate is $q/a$, and the lift through the \etale chart is uniquely
determined by Hensel's
lemma.  Write $d=[F'(Z):F']$, and let $P_a$ be the minimal polynomial
of $a$.  The horizontal divisor of $P_a(x)$ is $\overline Z$, with
multiplicity one.  Since all conjugates of $a$ have valuation $v(a)$,
the vertical order along the branch $x=0$ is
$v(N_{F'(a)/F'}(a))=d\,v(a)$; along the other branch it is zero.
Choose $a_c\in\Omega$ with $v(a_c)=v(a)$ and lift $(a_c,q/a_c)$ through
the chart to obtain $c_Z$.
The quotient
\[
                  \frac{P_a(x)}{(x-a_c)^d}
\]
is a local equation for $\overline Z-d\,\overline{c_Z}$.  On the
normalized special branches it reduces to $1$ and
$\overline{N_{F'(a)/F'}(a)/a_c^d}$, respectively, after canceling the
vertical orders.  Both are nonzero constants.  The associated line bundle is
therefore trivial on each normalized component, and in particular has
multidegree zero.  Being Cartier can be checked after the faithfully
flat completion; away from the node it is already clear.  The same
calculation applies after a constant extension, since the component and
the valuation specifying $c_Z$ are unchanged.

The sum of these corrected divisors is Cartier of multidegree zero,
and its generic class is
$P-[\sum m_Z[F'(Z):F']c_Z]$.  It defines a point of
$\mathcal G(\mathcal O_{F'})$, and hence the required point of
$\mathcal G(\mathcal O_E)$.
\end{proof}

\begin{lemma}[Preservation of Kummer classes]
\label{loc:kummer-transfer}
Retain the preceding notation and assume that the residue field of $E$
is separably closed.  For $u:T_pS_1\to T_pS_2$ as in
Proposition~\ref{loc:tate-realization}, $P\in S_1(E)$, and $n\ge1$, there
is $Q_n\in S_2(E)$ such that
\[
 u_n\kappa_{S_1,n}(P)=\kappa_{S_2,n}(Q_n)
       \quad\text{in }H^1(E,S_2[p^n]).
\]
\end{lemma}
\begin{proof}
Write $P=z+c$ as in Lemma~\ref{loc:point-correction}.  Since
$S_1(\Omega)$ is $p$-divisible, it is enough to treat the integral point
$z$.  Work over the strictly henselian ring $\mathcal O_E$.
The quotient of $\mathcal G_i[p^n]$ by its open subgroup
$\widehat{\mathcal G}_i[p^n]$ is \etale.
Smooth-group torsors over a strictly henselian local ring are trivial
\cite[Expos\'e~XXIV, Proposition~8.1 and its footnote]{SGA3}.  The fppf
cohomology sequences show that every $\mathcal G_i[p^n]$-torsor lifts
to a $\widehat{\mathcal G}_i[p^n]$-torsor and is itself the division
torsor of a point of $\mathcal G_i(\mathcal O_E)$.

Lift the division torsor of $z$ to the connected subgroup
$\widehat{\mathcal G}_1[p^n]$.  Push it out by $F_u[p^n]$ and then by
the inclusion into $\mathcal G_2[p^n]$.  The resulting torsor is the
division torsor of some $Q_n\in\mathcal G_2(\mathcal O_E)$.
On generic fibres, $F_u[p^n]$ is the restriction of $u_n$, so this
construction gives the asserted equality of Kummer classes.
\end{proof}

\subsection{The differential estimate}

The integral comparison between Kummer classes and differentials is the
step carried out by $p$-adic Hodge theory in
\cite[\S9, before Lemma~9.2]{Mochizuki1999}.  Here it follows from the
formal homomorphism and the point correction above: equality of finite
Kummer classes gives divisibility of the difference of the logarithmic
derivatives in the integral Lie algebra of the fixed model.

For a derivation $D$ of an $R_0$-algebra and a point $z$ of a smooth
commutative group, write $\partial_Dz$ for its logarithmic derivative,
characterized by
$\langle\eta,\partial_Dz\rangle=\langle z^*\eta,D\rangle$ for invariant
differentials $\eta$.  This operation is additive and is natural under
group homomorphisms.

\begin{lemma}[Integral differential comparison]
\label{loc:differential-comparison}
Let $u$ and $F_u$ be as in Proposition~\ref{loc:tate-realization}, and
let $F/K_0$ and $E/F$ be as in Lemma~\ref{loc:point-correction}.
Suppose $D:E\to E$ is a $K_0$-derivation preserving $\mathcal O_E$, with
$D(F)\subset F$ and continuous restriction to $F$.  For
$x\in\widehat{\mathcal G}_1(\mathfrak m_F)$ and every $n\ge1$,
\[
 \kappa_{S_2,n}(F_u(x))=u_n\kappa_{S_1,n}(x),\qquad
 \partial_D(F_u(x))=dF_u(\partial_Dx).
\]
The first equality is in $H^1(F,S_2[p^n])$.
If $S_2=J$ and $z\in J(E)$ has
$\kappa_{J,n}(z)=u_n\kappa_{S_1,n}(x)$, then
\[
 \partial_Dz-dF_u(\partial_Dx)
 \in p^n\operatorname{Lie}(\mathcal G_2)\otimes_{R_0}\mathcal O_E.
\]
\end{lemma}
\begin{proof}
The series $F_u$ converges at $x$, with
$F_u(x)\in\widehat{\mathcal G}_2(\mathfrak m_F)$.
Finite flatness of $[p^n]$ gives a division
point $y$ of $x$ in the formal neighborhood over a finite extension of
$F$.  For $\sigma\in G_F$ one has
\[
 \sigma F_u(y)-F_u(y)=F_u(\sigma y-y)=u_n(\sigma y-y).
\]
This proves the equality of Kummer classes.  Naturality of invariant
differentials, together with continuity of $D$, gives the equality of
logarithmic derivatives.

For the final assertion, the Kummer exact sequence gives
$z-F_u(x)=[p^n]w$ with $w\in J(E)$.  Apply
Lemma~\ref{loc:point-correction} to write $w=w_0+a$, where
$w_0\in\mathcal G_2(\mathcal O_E)$ and $a\in J(\Omega)$.
Since a $K_0$-derivation vanishes on $\Omega$,
\[
 \partial_Dz-dF_u(\partial_Dx)=p^n\partial_Dw_0.
\]
Finally, $\partial_Dw_0$ belongs to
$\operatorname{Lie}(\mathcal G_2)\otimes_{R_0}\mathcal O_E$ because $D$
preserves $\mathcal O_E$.  This proves the inclusion.
\end{proof}

\section{From finite Kummer classes to curve points}
\label{sec:approximation}

To pass from points of Jacobians to points of curves, we follow the
Picard-torsor argument of \cite[\S4, Proposition~7.4 and \S8]{Mochizuki1999}
and its presentation in \cite[\S5, Proposition~9 and Definition~10;
\S6]{Faltings1998}. A sufficiently divisible lift of a section
trivializes the degree-one Picard torsor. After passing to a cover of
suitable genus, the pro-$p$ hypothesis on the absolute Galois group
turns a degree-one Picard class into a rational point. We obtain points whose
sections agree with a prescribed section in a fixed finite Kummer quotient.

For a proper curve $C$ with a chosen geometric base point,
$\alpha_C:C\to J_C$ denotes the corresponding Abel--Jacobi morphism,
defined over the field of that point. For $\Gm$, the notation
$\alpha_{\Gm}$ denotes the identity. The differential pullback in the
proper case is independent of the base point, since translations on a
Jacobian preserve invariant differentials.

\begin{lemma}[A divided torsor]
\label{lem:divided-torsor}
Let $A/F$ be an abelian variety in characteristic zero, let $T$ be an
$A$-torsor, and let $s:G_F\to\Pi_T$ be a section. Let $m$ be a power of
$p$. Suppose
$E/F$ is finite of degree dividing $m$ and $P\in T(E)$ has the same $A[m]$-conjugacy class of sections as $s|_{G_E}$ in the
quotient whose geometric kernel is $A[m]$. Then $T(F)\ne\varnothing$.
\end{lemma}
\begin{proof}
The cocycle description of the extension with geometric kernel $A[m]$
associates to the splitting induced by $s$ a class
$\eta\in H^1(F,A)$ with $m\eta=[T]$
\cite[Proposition~2.2 and Lemma~2.3]{HarariSzamuely2009}.
In this description, agreement with the point section of $P$ means
that $\eta|_E=0$.
Restriction followed by corestriction gives $[E:F]\eta=0$. Since
$[E:F]\mid m$, it follows that $[T]=m\eta=0$.
\end{proof}

\begin{lemma}[Degree-one Picard classes over pro-$p$ fields]
\label{lem:picard-point}
Let $F$ have characteristic zero and pro-$p$ absolute Galois group. Let
$C/F$ be a smooth proper geometrically connected curve of genus $g$.
If $\Pic_C^1(F)\ne\varnothing$ and either $p$ is odd or $g$ is odd, then
$C(F)$ is nonempty. In fact its points avoid any prescribed finite closed
subscheme.
\end{lemma}
\begin{proof}
Let $\xi\in\Pic_C^1(F)$, and let $\beta\in\operatorname{Br}(F)$ be the
obstruction to representing $\xi$ by a line bundle over $F$
\cite[\S8.1, Proposition~4]{BLR1990}. Choose a
degree-one representative $M$ and descent isomorphisms over a finite Galois
extension $F'/F$, whose scalar cocycle $c_{\sigma,\tau}$ represents
$\beta$. This extension
has $p$-power degree, so restriction--corestriction shows that $\beta$ has
$p$-power order.

We use the determinant of cohomology to obtain an annihilator prime to
$p$; this is the Picard-obstruction argument recalled by Faltings
\cite[p.~143]{Faltings1998}.
Choose $m>2g-2$ with $p\nmid m(m+1-g)$, which is possible under the
stated parity hypothesis. Riemann--Roch gives
$h^0(C_{F'},M^{\otimes m})=m+1-g$. Taking determinants of the
semilinear maps on this space of sections makes the scalar cocycle
$c_{\sigma,\tau}^{m(m+1-g)}$ a coboundary. Hence
$m(m+1-g)\beta=0$, and $\beta=0$.
Thus $\xi$ is represented by a degree-one line bundle $L$ over $F$.

For a sufficiently large $d$ prime to $p$, the line bundle $L^{\otimes d}$
has a section whose zero divisor avoids the prescribed finite subscheme:
the evaluation map is surjective by Riemann--Roch, and $F$ is infinite.
The resulting effective divisor has degree $d$. Every closed point has
$p$-power degree, so at least one point in its support has degree one.
\end{proof}

We now apply these observations to a twist of an abelian cover of the
target. The use of covers carrying a lift of a section is the method
underlying Tamagawa's criterion for geometric sections
\cite[Corollary~2.10]{Tamagawa1997}. Here we require agreement with the
section only modulo a prescribed power of $p$; the pro-$p$ hypothesis
and the preceding two lemmas provide the required rational point.

\begin{lemma}[Finite point approximation]
\label{lem:point-approximation}
Let $X/K$ be a smooth proper geometrically connected curve or $\Gm$, let
$Y/K$ be proper hyperbolic, and let $\theta:\Pi_X\to\Pi_Y$ be over $G_K$.
Let $L\supset\bar K$ have pro-$p$ absolute Galois group. Suppose that on
every finite geometric $p$-cover induced by $\theta$, the corresponding
Tate map preserves finite Kummer images over every finite extension of
$L$. Then, for every $x\in X(L)$ and $n\ge1$, there is $y_n\in Y(L)$
whose point section agrees with $\theta s_x$ after geometric
abelianization modulo $p^n$.
\end{lemma}
\begin{proof}
Put $s=\theta s_x$ and choose $a\ge\max(n,2)$.
A point of $Y(\bar K)$ supplies a splitting, and
$G_L$ acts trivially on $T_pJ_Y/p^a$. Thus the abelianized extension
modulo $p^a$ is $(T_pJ_Y/p^a)\times G_L$. Write
$s(g)=(\rho_a(g),g)$. The quotient $(q,g)\mapsto q-\rho_a(g)$ defines a
geometrically connected $(T_pJ_Y/p^a)$-torsor $Y_a^s\to Y_L$ to which
$s$ lifts.
The Abel--Jacobi morphism from $Y_a^s$ to its degree-one Picard torsor
therefore carries the lifted section to a section of
$\Pic^1_{Y_a^s/L}$.

Let $T/L$ be the finite extension cut out by $\ker\rho_a$. Over $T$ the
twisted cover becomes a geometric cover $V$, and $x$ lifts to a point $z$
of the corresponding pullback cover $U$. Choose all covers and their
constant base points over a finite extension of $K$ inside $\bar K$.
Then $U$ is a proper curve or, in the multiplicative case, a copy of
$\Gm$ with the chosen base point corresponding to $1$.
These choices identify the Picard torsor with $J_V$ over $T$.
The arithmetic base-point cocycle vanishes on $G_T$, since $T$ contains
$\bar K$. Relative to the chosen constant point of $\Pic^1_V$, the
difference of sections at level $p^m$ is consequently the image of
$\kappa_{J_U,m}(\alpha_U(z))$ in the proper case, or
$\kappa_{\Gm,m}(z)$ in the multiplicative case, under the homomorphism
induced by $\theta$ on geometric abelianizations. This is the
section-difference calculation of \cite[Lemma~2.6]{Tamagawa1997},
applied to the finite $p$-primary quotient.

Take $m$ with $[T:L]\mid p^m$. By the Kummer-image hypothesis this
difference is the Kummer class of a point of $J_V(T)$. Translating the
constant Picard point by that point gives a point whose section agrees
with the prescribed section modulo $p^m$. Applying
Lemma~\ref{lem:divided-torsor} to the Picard torsor over $L$ gives
$\Pic^1_{Y_a^s/L}(L)\ne\varnothing$.
Riemann--Hurwitz gives
\[
 g(Y_a^s)-1=p^{2g(Y)a}(g(Y)-1).
\]
Thus its genus is odd when $p=2$, and Lemma~\ref{lem:picard-point}
produces an $L$-point of $Y_a^s$. Its image $y_n$ has the required
section difference modulo $p^a$, and hence modulo $p^n$.
\end{proof}

\section{Reconstruction from differentials}
\label{sec:reconstruction}

The reconstruction of a curve morphism by preserving relations among
regular differentials is the argument of Mochizuki
\cite[Proposition~12.1 and Corollary~13.4]{Mochizuki1999}; see also
Faltings \cite[\S6]{Faltings1998}.
The points just constructed need not be compatible as $n$ varies.
Their differential pullbacks nevertheless approach a fixed linear map,
and this will suffice to reconstruct a morphism. The reason is that a
proper hyperbolic curve is recovered from ratios of regular
differentials together with their first derivatives. The derivatives
retain the information lost by the canonical map in the hyperelliptic
case. A derivation bounded for the valuation then allows the algebraic
relations among these functions to pass to the limit.

\subsection{Canonical ratios and their first derivatives}

\begin{lemma}
\label{lem:first-jets}
Let $C/k$ be a proper hyperbolic curve in characteristic zero, with basis
$\omega_0,\ldots,\omega_{g-1}$ of its regular differentials. Set
\[
 U=\{\omega_0\ne0\},\qquad
 r_i=\omega_i/\omega_0,\qquad h_i=dr_i/\omega_0.
\]
Then $U$ is affine and
\begin{equation}
             \cO(U)=k[r_i]_i+\sum_i k[r_j]_j h_i=k[r_i,h_i]_i.
 \label{eq:jet-generation}
\end{equation}
Consequently dominant rational maps from an integral $k$-scheme to $C$
are determined by their pullbacks of regular differentials.
\end{lemma}
\begin{proof}
The canonical map is an embedding unless $C$ is geometrically
hyperelliptic \cite[Lemma~10.4(1)--(2)]{Mochizuki1999}.
In the hyperelliptic case, uniqueness of the involution makes it descend
to $k$. Its quotient is a smooth genus-zero curve, and the canonical
map factors through this quotient, followed by the embedding defined by
a complete linear system of degree $g-1$.
The inverse image of the affine canonical chart is $U$, so
$\cO(U)$ is finite over $k[r_i]_i$. In the nonhyperelliptic case,
$\cO(U)=k[r_i]_i$.

In the hyperelliptic case, the quotient has genus zero, so the involution
$\iota$ acts as $-1$ on regular differentials. Its invariant ring is $k[r_i]_i$, and
$\cO(U)=k[r_i]_i\oplus\cO(U)^-$, where the second summand is the
anti-invariant part. Invariant differentials on $U$ are precisely the
pullbacks of differentials from $\Spec k[r_i]_i$. Indeed, away from
the branch locus this follows from \etale descent; at a branch point,
the local equation $z^2=b$ gives $z\,dz=db/2$ and proves the same
assertion. Since $\omega_0$ trivializes the differentials on $U$, their
invariant part is $\cO(U)^-\omega_0$. The coordinate differentials
$dr_i=h_i\omega_0$ generate this module over $k[r_i]_i$. Thus
$\cO(U)^-=\sum_i k[r_j]_j h_i$, which proves
\eqref{eq:jet-generation}.

For the last assertion, two dominant maps with the same differential
pullback have the same values on every $r_i$. Differentiating and dividing
by the common nonzero pullback of $\omega_0$ gives the same values on
every $h_i$. Equation~\eqref{eq:jet-generation} proves equality of the
maps on the generic point.
\end{proof}

\begin{proposition}[Reconstruction of a morphism]
\label{prop:reconstruction}
Let $k\subset E\subset L$ be characteristic-zero fields with $E/k$
finitely generated of transcendence degree one. Let $v$ be a real-valued
nonarchimedean valuation on $L$, and let $D:L\to L$ be a derivation satisfying
\[
 D(k)=0,\quad D(E)\subset E,\quad D|_E\ne0,\qquad v(Db)\ge v(b).
\]
Let $C/k$ be proper hyperbolic and let
$a:H^0(C,\omega_C)\to\Omega^1_{E/k}$ be a nonzero $k$-linear map.
Suppose that for every $N$ there exists $y_N\in C(L)$ such that, on one
basis $\omega_i$,
\[
 v\bigl(\langle y_N^*\omega_i,D\rangle-
             \langle a(\omega_i),D\rangle\bigr)\ge N.
\]
Then there is a unique nonconstant $q\in C(E)$ with $q^*=a$.
\end{proposition}
\begin{proof}
Evaluation at $D$ identifies $\Omega^1_{E/k}$ with $E$. Put
$A_i=\langle a(\omega_i),D\rangle$ and reorder the basis so that
$A_0\ne0$. Write $\beta_{i,N}=\langle y_N^*\omega_i,D\rangle$.
For large $N$, $\beta_{0,N}\ne0$, and hence $y_N\in U(L)$, with $U$
as in Lemma~\ref{lem:first-jets}. Its evaluation map
$\varphi_N:\cO(U)\to L$
satisfies
\[
 \varphi_N(r_i)=\frac{\beta_{i,N}}{\beta_{0,N}},\qquad
 \varphi_N(h_i)=\frac{D\varphi_N(r_i)}{\beta_{0,N}}.
\]
The denominators have eventually constant valuation, and $D$ does not
decrease the valuation of an error. Therefore these values converge, in
the valuation sense, to
\[
 A_i/A_0,\qquad D(A_i/A_0)/A_0\quad\text{in }E.
\]
Every polynomial relation among the $r_i,h_i$ remains true at these
limits. Indeed, evaluation of a polynomial is continuous for the
valuation, and all the proposed limits already belong to $E$; no
completeness assumption on $E$ or $L$ is involved.
Equation~\eqref{eq:jet-generation}
therefore gives a $k$-algebra map $\psi:\cO(U)\to E$ and a point
$q\in U(E)$. For every $b\in\cO(U)$, both $\varphi_N(b)\to\psi(b)$ and
$D\varphi_N(b)\to D\psi(b)$.

Write each regular differential on $U$ as a finite sum
$\omega_i=\sum_j c_{ij}\,db_j$, where $b_j$ run through the generators
$r_i,h_i$ and $c_{ij}\in\cO(U)$. Passing to the limit in this expression
shows $\langle q^*\omega_i,D\rangle=A_i$. Thus $q^*=a$; in particular
$q$ is nonconstant. Uniqueness follows from Lemma~\ref{lem:first-jets}.
\end{proof}

\subsection{A valued generic point}

To apply Proposition~\ref{prop:reconstruction}, we realize the generic
point of the source curve in a valued field on which differentiation is
bounded. We also need the residue field to be separably closed and the
absolute Galois group to be pro-$p$, as required by the approximation
argument. The following construction provides these properties
simultaneously. The compositum of an unramified extension with the
algebraic constants is the valued field used by Faltings
\cite[\S6, pp.~146--147]{Faltings1998}. We record the boundedness of a
derivation on this field, which will allow us to pass algebraic
relations to the limit.

\begin{lemma}
\label{lem:gauss}
Let $K_0$ be a complete discretely valued field of mixed characteristic
$(0,p)$ and let $X/K_0$ be a smooth geometrically connected curve.
After a finite extension of $K_0$, its function field embeds in
$F=\widehat{K_0(t)}$, completed for the Gauss valuation, so that its
generic point lies in any prescribed formal neighborhood of a chosen
rational smooth point. On
\[
 L=F^{\nr}\Omega,\qquad \Omega=\overline{K_0},
\]
the derivation $D=\partial/\partial t$ satisfies
\begin{equation}
 D(\Omega)=0,\quad D(K_0(X))\subset K_0(X),\quad
 D|_{K_0(X)}\ne0,\quad v(Db)\ge v(b).
 \label{eq:bounded-D}
\end{equation}
The field $L$ and all its finite extensions have separably closed residue
fields, value group $\mathbf Q$ when $v(p)=1$, and pro-$p$ absolute
Galois groups.
\end{lemma}
\begin{proof}
Choose a rational smooth point after finite extension, an \etale coordinate
$z$ vanishing there, and a standard \etale equation $P(z,y)=0$ normalized
by $P_y(0,0)=1$. With a uniformizer $\pi$ and sufficiently large $m$,
\[
 \pi^{-m}P(\pi^{2m}t,\pi^m w)=w+\pi H(t,w),
                 \qquad H\in\cO_{K_0}[t,w].
\]
Hensel's lemma gives a point over $F$ with $z=\pi^{2m}t$. It is generic,
and $m$ can be increased to place it in any prescribed formal
neighborhood. Differentiation does not decrease Gauss valuation on
$K_0(t)$, and consequently extends continuously to $F$. Separability
implies that it preserves the embedded function field.

For each finite constant extension $K_1/K_0$, the field $FK_1$ is the
Gauss completion of $K_1(t)$, so the same estimate holds. Formal
\'etaleness gives unique extension of a derivation across a finite
\etale algebra \cite[Tag~\stacks{00UR}]{Stacks}. Applying this to the unramified
extensions, after constant extension, shows that $D$ preserves
$\cO_L$. Every value in $L$ is realized by a constant in $\Omega$.
Dividing an element by a constant of the same value proves the full
inequality in \eqref{eq:bounded-D}.

The residue field of $L$ is separably closed. Since $L/F$ is algebraic
and $\Omega$ contains all roots of $p$, its value group is $\mathbf Q$;
these statements persist under finite extension. For every prime
$\ell\ne p$, each such extension $T$ has $\ell$-divisible multiplicative
group: divide the value by $\ell$, then apply Hensel's lemma to a unit.
Kummer theory excludes cyclic extensions of degree $\ell$ over every
finite extension of $L$. A non-$p$ finite Galois quotient would contain
a subgroup of order $\ell$ and yield just such an extension. Hence $G_L$
is pro-$p$.
\end{proof}

\subsection{The geometric curve morphism}

Let $S_1,S_2$ be Jacobians or $\Gm$ over a $p$-subtame field $K$, and
let $u:T_pS_1\to T_pS_2$ be $G_K$-equivariant. Fix an embedding
$K\subset M^{\tame}\subset\Omega$ as in Section~\ref{sec:fields}.
Lemma~\ref{lem:tate-descent} and Proposition~\ref{loc:tate-realization}
define an $\Omega$-linear differential
\begin{equation}
 \lambda_u:\Lie(S_1)_\Omega\longrightarrow\Lie(S_2)_\Omega.
 \label{eq:lambda}
\end{equation}
It is compatible with extension of the finite field of definition,
composition, and algebraic group morphisms. The proposition below
constructs a geometric curve morphism from this differential. Its
descent to $K$, and the comparison with the original nonabelian
homomorphism, will be proved in the next section.

\begin{proposition}
\label{prop:geometric-map}
Let $K$ be $p$-subtame, $X/K$ a smooth proper geometrically connected
curve or $\Gm$, and $Y/K$ proper hyperbolic. Choose geometric base
points on the proper curves, and use $1$ on $\Gm$.
Let $\theta:\Pi_X\to\Pi_Y$ be over $G_K$.
Its geometric abelianization $\theta^{\ab}$ has target $T_pJ_Y$ and
source $T_pJ_X$ in the proper case, or $\Zp(1)$ when $X=\Gm$.
If
\[
 a_\theta=\alpha_X^*\lambda_{\theta^{\ab}}^*:
 H^0(Y_\Omega,\omega_{Y_\Omega})\longrightarrow
 H^0(X_\Omega,\Omega^1_{X_\Omega/\Omega})
\]
is nonzero, there is a unique nonconstant morphism
$f_\Omega:X_\Omega\to Y_\Omega$ with $f_\Omega^*=a_\theta$.
The construction commutes with morphisms of the relevant finite covers
and their geometric deck transformations.
\end{proposition}
\begin{proof}
Descend the curves, base points, and $\theta^{\ab}$ to a finite extension $M_0/M$
of the local field using Lemma~\ref{lem:tate-descent}. Enlarge $M_0$ to choose the
semistable models and the integral formal map $F_{\theta^{\ab}}$. Apply
Lemma~\ref{lem:gauss} near a base point of $X$ so that its generic point
$x\in X(F)$ has Abel--Jacobi image $\alpha_X(x)$ in the formal
neighborhood of the identity. Set
$L=F^{\nr}\Omega$ and
$q=F_{\theta^{\ab}}(\alpha_X(x))\in J_Y(F)$.
Lemma~\ref{loc:differential-comparison} identifies the Kummer class of
$q$ modulo $p^n$ with the image under $(\theta^{\ab})_n$ of the Kummer
class of $\alpha_X(x)$, taken in $J_X$ or $\Gm$ according to the
source. It also gives
\[
 \partial_Dq=\lambda_{\theta^{\ab}}(\partial_D\alpha_X(x)).
\]

To obtain points of $Y$ approximating $q$, we check the hypothesis of
Lemma~\ref{lem:point-approximation} over this same field $L$.
The untwisted covers $V,U$ occurring there and their induced Tate maps
descend to a finite extension of $K$.
Lemma~\ref{lem:tate-descent} puts these data and their models over some
finite $M_1/M_0$ inside $\Omega$. For any finite $T/L$,
the field $FM_1\subset T$ is a complete discretely valued field and
$T/FM_1$ is algebraic with separably closed residue field. Thus
Lemma~\ref{loc:kummer-transfer} preserves every finite Kummer image over
$T$ itself. Finite geometric $p$-covers of $\Gm$ become
$t\mapsto t^{p^b}$ after finite constant extension
\cite[Example~9.8]{Milne1986}, so the same argument includes that case.

Lemma~\ref{lem:point-approximation} now supplies $y_n\in Y(L)$ whose
Abel--Jacobi images have the same Kummer classes as $q$ modulo $p^n$.
Apply Lemma~\ref{loc:differential-comparison}, using the one fixed lattice
of invariant differentials on the model of $J_Y$. Their pullbacks to $Y$
form a basis by \eqref{eq:AJ}. We obtain the approximation required in
Proposition~\ref{prop:reconstruction}, with
$M_0\subset M_0(X)\subset L$. This gives a nonconstant rational map over
$M_0$, which extends because the source is a smooth curve and the target
is proper \cite[Tag~\stacks{0BXZ}]{Stacks}. Base change gives $f_\Omega$.

Uniqueness follows from Lemma~\ref{lem:first-jets}. For a comparison
between covers, all the data can be defined over a common finite
extension of the local field. Naturality of
$\lambda_{\theta^{\ab}}$, followed by the same uniqueness, proves the assertions
about finite covers and geometric deck maps.
\end{proof}

\section{The Hom theorem for proper curves}
\label{sec:proper-hom}

The preceding section constructs a geometric morphism with the
differential prescribed by $\theta$. We must still prove that it induces
$\theta$ on the fundamental group. We follow Mochizuki's comparison
over finite covers \cite[\S14, proof of Theorem~14.1]{Mochizuki1999},
also used by Faltings \cite[proof of Theorem~14]{Faltings1998}.
The construction is applied to every finite cover of the target.
Its uniqueness makes the resulting
morphisms compatible and equivariant under deck transformations, so
that they recover the full geometric homomorphism. The following lemma
then gives descent to the ground field and equality on the arithmetic
fundamental groups. We include the descent and centralizer arguments
for the present ground fields.

\begin{lemma}[Geometric faithfulness and arithmetic descent]
\label{lem:arithmetic-descent}
Let $X/k$ be a smooth proper geometrically connected curve and $Y/k$
proper hyperbolic, in characteristic zero. A nonconstant morphism
$X\to Y$ is determined by its geometric outer pro-$p$ homomorphism.
If a nonconstant morphism over an algebraically closed extension $\Omega/k$
has the same geometric outer homomorphism as an open
$\theta:\Pi_X\to\Pi_Y$ over $G_k$, it descends to a unique $k$-morphism,
whose entire arithmetic homomorphism is $\theta$ modulo geometric
conjugation.
\end{lemma}
\begin{proof}
The Tate-module functor is faithful on abelian varieties
\cite[Lemma~12.2]{MilneAV1986}. Thus equal geometric outer
homomorphisms give equal Jacobian homomorphisms and equal differential
pullbacks, by the Abel--Jacobi identifications
\cite[Propositions~2.2 and 9.1]{Milne1986}.
Lemma~\ref{lem:first-jets} proves the first assertion.

For every $\sigma\in\Aut(\Omega/k)$, the morphisms $f_\Omega$ and
$\sigma f_\Omega$ have the same geometric outer homomorphism, by the
$G_k$-compatibility of $\theta$ and invariance of geometric fundamental
groups under algebraically closed extension. They are therefore equal.
Their function-field map descends to $k$. One may see the descent
directly from
\[
                         \Omega(X)^{\Aut(\Omega/k)}=k(X).
\]
To verify this equality, take an invariant fraction with a denominator
$\sum_{i=1}^r a_ie_i$, where the $e_i\in k(X)$ are $k$-linearly
independent, $a_1=1$, and $r$ is minimal. Subtracting any conjugate of
the denominator gives a denominator with fewer terms unless it is zero.
All its coefficients, and then those of the numerator, are fixed and
belong to $k$. The descended rational map extends by properness.

It remains to compare the arithmetic homomorphisms. We adapt the
deck-transformation argument of \cite[Lemma~1 and Corollary~2]{Faltings1998}
to the common open geometric image. After geometric
conjugation, assume $f_*$ and $\theta$ agree on
$\Delta_X$, whose common image is open in $\Delta_Y$.
For $\gamma\in\Pi_X$, the element
$z=f_*(\gamma)^{-1}\theta(\gamma)\in\Delta_Y$ centralizes
$\theta(\Delta_X)$.
If $z\ne1$, choose open normal $N\triangleleft\Delta_Y$ with
$N\subset\theta(\Delta_X)$ and $z\notin N$. On the corresponding proper hyperbolic
cover $C\to Y_{\bar k}$, the nontrivial deck transformation defined by
$z$ induces the identity on $N^{\ab}=T_pJ_C$. The faithfulness just
proved makes this deck transformation the identity, a contradiction.
Thus $f_*=\theta$ on the entire arithmetic group.
\end{proof}

\begin{proposition}
\label{prop:proper-curves}
The pro-$p$ assertion of Theorem~\ref{thm:main} holds when $X$ and $Y$
are smooth proper geometrically connected curves and $Y$ is hyperbolic.
\end{proposition}
\begin{proof}
Replace the target by the cover corresponding to the image of an open
$\theta$, so that $\theta$ is surjective. The induced homomorphism
$T_pJ_X\to T_pJ_Y$ is surjective, so the differential
$a_\theta$ of Proposition~\ref{prop:geometric-map} is nonzero. Fix the
resulting $f_\Omega:X_\Omega\to Y_\Omega$.

Let $V\to Y$ be an arithmetic Galois cover corresponding to an open
normal subgroup of $\Pi_Y$, and let $W\to X$ be its $\theta$-pullback.
Both are geometrically connected over the same finite constant field
$E/K$. The induced map $\Pi_W\to\Pi_V$ is again surjective.
Proposition~\ref{prop:geometric-map} gives $r:W_\Omega\to V_\Omega$.
Naturality of its differential and Lemma~\ref{lem:first-jets} imply
that $r$ lies over $f_\Omega$.
They also imply equivariance under every geometric deck transformation.
The resulting equivariant morphism
$W_\Omega\to X_\Omega\times_{Y_\Omega}V_\Omega$ is a morphism
between torsors under the same finite group, and is therefore an
isomorphism. The geometric quotients so obtained are cofinal:
every open subgroup of $\Delta_Y$ contains the intersection with
$\Delta_Y$ of some open normal subgroup of $\Pi_Y$.
Lemma~\ref{lem:cover-tools}(i) identifies the geometric outer
homomorphism of $f_\Omega$ with that of $\theta$.
Lemma~\ref{lem:arithmetic-descent} now supplies the required $K$-morphism
and equality of the full arithmetic homomorphisms. It also proves
uniqueness.

A nonconstant morphism of proper curves is finite. Its function-field
map has open Galois image, and the maps from the generic Galois groups
to the respective fundamental groups are surjective
\cite[Expos\'e~V, Proposition~8.2]{SGA1}. Hence its induced homomorphism is open.
\end{proof}

To allow an open source, it suffices to show that a homomorphism into
the fundamental group of a proper hyperbolic curve kills the inertia
at each missing point. The necessary local assertion is the following
consequence of reconstruction for $\Gm$.

\begin{lemma}[Vanishing of puncture inertia]
\label{lem:puncture}
For $K$ $p$-subtame and $Y/K$ proper hyperbolic, every
$G_K$-homomorphism $\Pi_{\Gm}\to\Pi_Y$ kills $\Delta_{\Gm}$.
Consequently every $\Pi_U\to\Pi_Y$, for a smooth separated curve $U$, factors
through the smooth proper compactification of $U$, with unchanged image.
\end{lemma}
\begin{proof}
Write $\Pi_{\Gm}=\Zp(1)\rtimes G_K$, using the point $1$.
Suppose a generator has nontrivial image $z\in\Delta_Y$. Detect $z$
in a finite arithmetic quotient and take the inverse image in
$\Delta_Y$ of the cyclic subgroup generated by its image. This open
subgroup is stable under the image section of $G_K$, since conjugation
sends $z$ to $z^{\chi_p(g)}$. It gives a proper hyperbolic cover $C\to Y$
to which the homomorphism lifts. The induced map
$\Zp(1)\to T_pJ_C$ is nonzero, as is its differential by
Proposition~\ref{loc:tate-realization}. Proposition~\ref{prop:geometric-map}
would produce a nonconstant $\mathbf G_{m,\Omega}\to C_\Omega$. Properness extends
it to $\PP^1_\Omega$, contradicting Riemann--Hurwitz.

For a boundary point of $U$, first make the point rational over a finite
extension $K'/K$. A parameter and a compatible system of its roots give
a splitting of the decomposition group of the henselian local field
as $\widehat{\mathbf Z}(1)\rtimes G_{K'}$.
The action on inertia is cyclotomic. Since the geometric kernel of
$\Pi_Y$ is pro-$p$, the resulting homomorphism factors through
$\Zp(1)\rtimes G_{K'}$. The first assertion kills this inertia.
A finite cover extends \etale over the compactification precisely when
these inertia groups act trivially: this is the normalization
description of covers of a normal scheme
\cite[Expos\'e~I, Corollaire~10.3]{SGA1}.
Thus their closed normal closure is the kernel of the map to the
fundamental group of the compactification. This proves the
factorization; passage to the inverse limit over the open subcurves
gives the same assertion for the generic point.
\end{proof}

Combining Proposition~\ref{prop:proper-curves} with
Lemma~\ref{lem:puncture} proves the pro-$p$ Hom theorem from every smooth
separated curve to a proper hyperbolic curve. For a nonseparated smooth
curve, apply this result on affine opens and glue by uniqueness.

\section{Normal sources and open targets}
\label{sec:global}

We next remove the restrictions on the source and target. For the source,
the essential point is to find a curve over a larger ground field which
preserves every connected geometric cover of $X$. The curve theorem
then gives a rational map on $X$. A separate extension argument shows
that its compatibility with the fundamental group excludes
indeterminacy, even when $X$ is singular or nonseparated. For an open
target, ramification at the missing points will similarly exclude
intersection with the boundary.

\subsection{A simultaneous curve reduction}

An arbitrary curve in $X$ need not preserve the image of $\theta$ on
geometric fundamental groups. We instead use a generic pencil in the
function field, with independent parameters in the ground field. Its
regularity after every finite extension of the function field is what
ensures the required surjectivity. Mochizuki obtains this surjectivity
from a Bertini pencil and the Lefschetz theorem
\cite[Lemmas~16.1--16.2]{Mochizuki1999}; the following field-theoretic
argument supplies the pencil without choosing a smooth projective
model of $X$.

\begin{lemma}
\label{lem:curve-reduction}
Let $F/k$ be a regular finitely generated extension of positive
transcendence degree $d$, in characteristic zero. There are algebraically
independent variables $\mathbf t$ over $F$ and a purely transcendental
extension $E/k$ contained in $F(\mathbf t)$ such that
$F(\mathbf t)/E$ is regular of transcendence degree one. They can be
chosen so that, for every finite $k'/k$ and every finite extension
$M/Fk'$ regular over $k'$, the extension $M(\mathbf t)/Ek'$ is regular.
If $F=k(X)$ for a normal geometrically integral finite-type scheme,
there is a smooth geometrically connected curve $U/E$ with function
field $F(\mathbf t)$ and a $k$-morphism $U\to X$ inducing a surjection
$\Delta_{U/E}\to\Delta_{X/k}$.
\end{lemma}
\begin{proof}
We first prove the field-theoretic assertion underlying the construction.
If $M/k$ is
regular and finitely generated and $a,b\in M$ are algebraically
independent, then
$M(t)/k(t,at+b)$ is regular. Let $N_0$ be the relative algebraic closure
of $k(a,b)$ in $M$, which is finite over $k(a,b)$.
For $z\in M(t)$ algebraic over $k(t,at+b)$, one has $z\in N_0(t)$.
The intermediate field $k(a,b,t,z)$ is $L_0(t)$ for some
$k(a,b)\subset L_0\subset N_0$. Indeed, take a primitive element of
$N_0/k(a,b)$. The coefficients of its minimal polynomial over
$k(a,b,t,z)$ are algebraic over $k(a,b)$, hence belong to $N_0$; the
subfield $L_0$ they generate has the asserted property by comparison
of degrees.
Thus $k(a,b,t,z)$ is stable under $\sigma:t\mapsto t+1$ fixing $M$.
Put $u=at+b$ and $u'=u+a$. Then
$k(a,b,t,z)=k(t,u,z)(u')$ is regular over $k(t,u')$, because
$k(t,u,z)/k(t)$ is regular and $u'$ is independent of it.
But $\sigma z$ is algebraic over $k(t,u')$, so
$\sigma z\in k(t,u')$, giving $z\in k(t,u)$.

Now choose a transcendence basis $x_1,\ldots,x_d$ of $F/k$ and put
\[
 u_i=x_i+t_i x_d\quad(1\le i<d),\qquad
 E=k(t_i,u_i:1\le i<d).
\]
The reversible change $x_i=u_i-t_i x_d$ gives
$\trdeg_E F(\mathbf t)=1$. Apply the pencil observation successively,
also to every $M/k'$ under consideration. It proves the asserted
regularity with the same choices for all finite covers. Adding
independent variables preserves their degrees and Galois groups.

It remains to realize this field construction geometrically. Choose a
nonempty affine open of $X$. Clearing denominators of its
finitely many coordinate functions in $F(\mathbf t)$ realizes a
nonempty smooth affine model $U/E$ with a map to this open.
Every connected geometric finite cover of $X$ descends to some finite
$k'/k$; its function field is one of the $M$ above. The proved
regularity keeps its pullback to $U_{\bar E}$ connected. This is exactly
the asserted surjectivity on geometric fundamental groups, and hence
on their pro-$p$ quotients.
\end{proof}

The next lemma supplies the extension across the omitted locus of $X$.
For smooth sources the extension is unconditional
\cite[Lemma~16.3]{Mochizuki1999}, by the classical extension theorem
for maps to abelian varieties. For a normal source we use the stated
factorization on fundamental groups. The proof uses only abelian
covers of the target. A positive-dimensional
fibre in the closure of the graph would contain the whole target curve;
the assumed factorization on fundamental groups would then trivialize
covers of arbitrarily large degree over one fixed finite algebra.

\begin{lemma}[Extension across a normal source]
\label{lem:normal-extension}
Let $X/k$ be normal, geometrically integral and of finite type in
characteristic zero, let $F=k(X)$, and let $C/k$ be a smooth proper geometrically connected curve of
genus $g\ge1$. A rational map $q:\Spec F\to C$ extends uniquely to
$X$ if its geometric pro-$p$ abelianization factors as
\[
 \Delta_{\Spec F/k}\longrightarrow\Delta_{X/k}
                         \longrightarrow T_pJ_C.
\]
\end{lemma}
\begin{proof}
Take the closure $\Gamma\subset X\times C$ of the graph. It is proper
birational over $X$. Suppose it has a positive-dimensional geometric
fibre over $\bar x=\Spec\Omega\to X$. Its reduced fibre is then the
whole $C_\Omega$. Choose an affine neighborhood $\Spec A$ of its generic
point in $\Gamma$ and let $B$ be the normalization of $A$ in $F$.
This is finite since $A$ is of finite type over a field. The algebra
\[
             T=B\otimes_A\Omega(C)
\]
is nonzero and finite-dimensional over $\Omega(C)$.

Choose $n$ with $p^{2gn}>\dim_{\Omega(C)}T$.
After a finite constant extension $k'/k$,
Lemma~\ref{lem:quotient-descent} realizes the quotient
$T_pJ_C/p^n$ as a geometrically connected abelian cover $V\to C_{k'}$
of degree $p^{2gn}$. The factorization assumption makes $q^*V$ the
generic fibre of a finite \etale torsor over $X_{k'}$: the kernel of
$G_{Fk'}\to\Pi_{X_{k'}}$ is geometric and is killed by the assumed
factorization. Over the normal
scheme $\Spec(B\otimes_k k')$, its pullback and the pullback of $V$
are isomorphic: finite \etale covers of a normal integral scheme are
determined by their generic fibres
\cite[Expos\'e~I, Corollaires~10.2--10.3]{SGA1}.
A section of that torsor at $\bar x$, pulled back to $T$, therefore gives
an $\Omega(C)$-algebra map $\Omega(V)\to T$. It is injective, since
$\Omega(V)$ is a field. This forces
$p^{2gn}\le\dim_{\Omega(C)}T$, a contradiction.

Thus $\Gamma\to X$ is proper with finite fibres, hence finite
\cite[Expos\'e~VIII, Th\'eor\`eme~6.1]{SGA1}.
Being birational, it is an isomorphism by normality of $X$.
This proves existence;
separatedness of $C$ proves uniqueness. 
\end{proof}

\begin{proposition}
\label{prop:proper-normal}
The pro-$p$ Hom theorem holds for normal geometrically integral
finite-type sources and proper hyperbolic targets over $p$-subtame
fields.
\end{proposition}
\begin{proof}
In dimension zero both sides are empty. Suppose that $\dim X>0$,
let $\theta:\Pi_X\to\Pi_Y$ be open over $G_K$, and put $F=K(X)$.
Apply Lemma~\ref{lem:curve-reduction} to obtain $U/E\to X$ with
$E(U)=F(\mathbf t)$. The field $E$ is $p$-subtame by
Lemma~\ref{lem:fields}. The induced map
\[
 \Pi_{U/E}\longrightarrow\Pi_{Y/K}\times_{G_K}G_E=\Pi_{Y_E/E}
\]
has the same geometric image as $\theta$ and is open. The curve case
therefore gives $q\in Y(F(\mathbf t))$ inducing this homomorphism.
For a proper curve of positive genus,
\[
                         Y(F(\mathbf t))=Y(F):
\]
remove the variables one at a time, since a rational map
$\PP^1_F\dashrightarrow Y$ extends and is constant by
Riemann--Hurwitz. The surjection $G_{F(\mathbf t)}\to G_F$ shows that
$q_*$ agrees with $\theta$ on the entire generic arithmetic group.
Lemma~\ref{lem:normal-extension} extends $q$ to $X$. Surjectivity of
$G_F\to\pi_1(X)$ \cite[Expos\'e~V, Proposition~8.2]{SGA1}
then gives the prescribed arithmetic homomorphism.
Uniqueness follows by pulling two candidates back to $U$, applying the
curve case, and descending equality from the generic point.
\end{proof}

\subsection{Avoiding the boundary of the target}

To treat an open target, we use proper compactifications of its finite
covers. We recall two covering facts about a hyperbolic curve $Y$ in
characteristic zero: each geometric pro-$p$ cusp inertia is a copy of
$\Zp(1)$, and $Y$ has an arithmetic finite $p$-cover whose smooth
compactified geometric components have genus at least two. Here
``$p$-cover'' refers to its geometric degree, not its constant-field
degree. When $Y$ is proper the second assertion is immediate. For an
open $Y$, both assertions follow from the usual presentation of its
geometric fundamental group
\cite[Expos\'e~XIII, Corollaire~2.12]{SGA1}. If there are $r\ge2$
boundary points, abelianizing this presentation gives residue vectors
whose coordinates sum to zero. The vectors
$(1,-1,0,\ldots,0)$ modulo $p^a$ detect cusp inertia to arbitrary
order. The cyclic cover with residues
$(1,\ldots,1,-(r-1))$ has, for $a>v_p(r-1)$,
\[
 2g(\bar V)-2
 =p^a(2g(\bar Y)-2+r)-(r-1)-p^{v_p(r-1)}>0
\]
when $a$ is large. If $r=1$, first take an unramified cover of degree
$p$ of $\bar Y$, which exists since $g(\bar Y)\ge1$, and then apply
the preceding construction. Finally descend the cover after finite
constant extension and take an arithmetic Galois closure using
Lemma~\ref{lem:quotient-descent}.

The maps to these proper compactifications are supplied by
Proposition~\ref{prop:proper-normal}. They descend to a map to
$\bar Y$. We then use Mochizuki's ramification argument
\cite[proof of Corollary~14.2]{Mochizuki1999}:
if that map met a missing point of $Y$, its order of
vanishing there would have to be divisible by the ramification indices
in every finite cover. The unboundedness of cusp inertia makes this
impossible.

\begin{theorem}[The pro-$p$ Hom theorem]
\label{thm:pro-p}
Theorem~\ref{thm:main} holds for $\Sigma=\{p\}$.
\end{theorem}
\begin{proof}
As before, replacing $Y$ by the image cover reduces the proof to a
surjective $\theta:\Pi_X\to\Pi_Y$. Choose an arithmetic Galois
cover $V_0\to Y$ as above, corresponding to $N_0\triangleleft\Pi_Y$.
For each open normal $N\subset N_0$, let $V\to Y$ be its cover, let
$W\to X$ be the $\theta$-pullback, and let $E$ be their common constant
field. Apply Proposition~\ref{prop:proper-normal} to the induced open
map from $\Pi_W$ to the fundamental group of the compactification
$\bar V$. It gives a unique dominant $E$-morphism $r_V:W\to\bar V$.
Uniqueness gives compatibility with $V\to V_0$ and with all arithmetic
deck transformations, including their semilinear action on $E$.
Hence $r_{V_0}$ descends to a dominant $\bar f:X\to\bar Y$.

Suppose $\bar f$ meets the boundary. Its inverse image is a nonempty
effective Cartier divisor, so has a height-one point $x$.
The local ring $\cO_{X,x}$ is a DVR. If $t$ is a parameter at the
boundary point $y$, put $e=\ord_x(\bar f^*t)>0$. For a cusp inertia
generator $z$, choose $N\subset N_0$ with $z^e\notin N$. The order
$r$ of $z$ in $\Pi_Y/N$ does not divide $e$, and every branch of
$\bar V$ above $y$ has ramification index $r$. At a point $w\in W$
above $x$, the cover $W\to X$ is \etale. A parameter $s$ at
$r_V(w)$ consequently gives
\[
              e=\ord_w(\bar f^*t)=r\,\ord_w(r_V^*s),
\]
a contradiction. Thus $\bar f$ factors through $f:X\to Y$.

All $r_V$ now land in $V$, and their arithmetic equivariance yields
$W\simeq X\times_YV$ as torsors, with identity on the constant
quotients. Lemma~\ref{lem:torsor-comparison} identifies $f_*$ with the
specified $\theta$ by one geometric conjugation. Uniqueness follows
after pulling back to $W_0$ and using uniqueness for the proper target
$\bar V_0$.

Finally, every dominant morphism $X\to Y$ induces an open
homomorphism. The relative algebraic closure of $K(Y)$ in $K(X)$ is
finite, and $K(X)$ is regular over it. Thus the generic Galois image is
open. Surjectivity from generic Galois groups to fundamental groups of
normal schemes \cite[Expos\'e~V, Proposition~8.2]{SGA1}
proves the assertion for $\Pi_X\to\Pi_Y$.
\end{proof}

\section{Passage to a set of primes}
\label{sec:sigma}

The passage from the pro-$p$ theorem to a set of primes is due to
Mochizuki. It appears for the full profinite statement in
\cite[Remark after Theorem~16.5]{Mochizuki1999}, and for arbitrary
$\Sigma$ containing $p$ in
\cite[Theorem~A and the following remark, p.~121]{Mochizuki2003}.
The same argument applies here because finite extensions of a
$p$-subtame field are again $p$-subtame. One applies the pro-$p$ theorem
to each finite $\Sigma$-cover and uses uniqueness to recover the cover
itself. We include the argument to make explicit its compatibility
with the constant-field quotient.

\begin{proof}[Proof of Theorem~\ref{thm:main}]
Let $\theta:\Pi_X^\Sigma\to\Pi_Y^\Sigma$ be open over $G_K$.
Lemma~\ref{lem:cover-tools} reduces it to a surjection. Passing to
the canonical relative pro-$p$ quotients gives a surjective
$\theta_p:\Pi_X\to\Pi_Y$. Theorem~\ref{thm:pro-p} gives a unique
dominant $f:X\to Y$ inducing $\theta_p$.

For every open normal $N\triangleleft\Pi_Y^\Sigma$, let $V\to Y$
and $W\to X$ be the corresponding torsor and its $\theta$-pullback.
They have the same
finite constant field $E/K$. By Lemma~\ref{lem:cover-tools}(ii), the
restricted map is a surjection
\[
                   \Pi_{W/E}^\Sigma\longrightarrow\Pi_{V/E}^\Sigma.
\]
Pass to its own relative pro-$p$ quotients and apply
Theorem~\ref{thm:pro-p} over $E$, which is $p$-subtame. This gives
a unique dominant $E$-morphism $r_V:W\to V$.

The composites $W\xrightarrow{r_V}V\to Y_E$ and
$W\to X_E\xrightarrow{f_E}Y_E$ induce the same relative pro-$p$
homomorphism over $E$. Pro-$p$ uniqueness therefore makes them equal.
For $a\in\Pi_Y^\Sigma/N$, the map $a_V^{-1}r_Va_W$ is $E$-linear, even if
$a_V$ and $a_W$ are semilinear over $E$. It induces the same pro-$p$
homomorphism as $r_V$, because the restricted homomorphism is
equivariant under conjugation by $a$. The same uniqueness gives
\[
                         r_Va_W=a_Vr_V.
\]
Thus the induced morphism is an isomorphism of $(\Pi_Y^\Sigma/N)$-torsors
$W\simeq X\times_YV$. Its action on the constant quotient is the
identity, since $r_V$ is over $E$. Lemma~\ref{lem:torsor-comparison}
now proves $f_*=\theta$ modulo one inner automorphism from
$\Delta_Y^\Sigma$.

Two morphisms inducing the same relative pro-$\Sigma$ homomorphism
modulo geometric conjugation induce the same such pro-$p$ homomorphism, so are
equal by Theorem~\ref{thm:pro-p}. Openness for dominant morphisms is
the same generic-Galois argument used there. This proves
\eqref{eq:main} for every $\Sigma$ containing $p$.
\end{proof}

\begin{thebibliography}{99}
\small
\raggedright
\hypersetup{urlcolor=black}

\bibitem{BLR1990}
S.~Bosch, W.~L\"utkebohmert, and M.~Raynaud,
\href{https://doi.org/10.1007/978-3-642-51438-8}{\emph{N\'eron Models}},
Ergebnisse der Mathematik und ihrer Grenzgebiete (3), vol.~21,
Springer-Verlag, Berlin, 1990.

\bibitem{DeligneMumford1969}
P.~Deligne and D.~Mumford,
\href{https://doi.org/10.1007/BF02684599}{\emph{The irreducibility of the
space of curves of given genus}},
Inst. Hautes \'Etudes Sci. Publ. Math. \textbf{36} (1969), 75--109.

\bibitem{SGA3}
M.~Demazure and A.~Grothendieck (dir.),
\href{https://doi.org/10.1007/BFb0059027}{\emph{Sch\'emas en groupes~III:
Structure des sch\'emas en groupes r\'eductifs (SGA~3)}},
S\'eminaire de g\'eom\'etrie alg\'ebrique du Bois Marie 1962--1964,
Lecture Notes in Mathematics, vol.~153,
Springer-Verlag, Berlin, 1970.

\bibitem{Faltings1998}
G.~Faltings, \emph{Curves and their fundamental groups
[following Grothendieck, Tamagawa and Mochizuki]},
S\'eminaire Bourbaki, Vol.~1997/98, Exp.~840,
Ast\'erisque \textbf{252} (1998), 131--150.

\bibitem{SGA1}
A.~Grothendieck (dir.),
\emph{Rev\^etements \'etales et groupe fondamental (SGA~1)},
S\'eminaire de g\'eom\'etrie alg\'ebrique du Bois Marie 1960--1961,
Lecture Notes in Mathematics, vol.~224, Springer-Verlag, Berlin, 1971;
revised edition, Documents Math\'ematiques, vol.~3,
Soci\'et\'e Math\'ematique de France, Paris, 2003.

\bibitem{Grothendieck1997}
A.~Grothendieck,
\href{https://doi.org/10.1017/CBO9780511758874.004}{\emph{Brief an
G.~Faltings}} (27 June 1983),
in \emph{Geometric Galois Actions~1: Around Grothendieck's
Esquisse d'un Programme}, L.~Schneps and P.~Lochak (eds.),
London Mathematical Society Lecture Note Series, vol.~242,
Cambridge University Press, Cambridge, 1997, pp.~49--58;
English translation, pp.~285--293.

\bibitem{HarariSzamuely2009}
D.~Harari and T.~Szamuely,
\href{https://doi.org/10.1007/s00208-008-0327-z}{\emph{Galois sections
for abelianized fundamental groups}},
with an appendix by E.~V.~Flynn,
Math. Ann. \textbf{344} (2009), no.~4, 779--800.

\bibitem{Hoshi2017}
Y.~Hoshi,
\href{https://doi.org/10.2206/kyushujm.71.1}{\emph{On the Grothendieck
conjecture for affine hyperbolic curves over Kummer-faithful fields}},
Kyushu J. Math. \textbf{71} (2017), no.~1, 1--29.

\bibitem{MilneAV1986}
J.~S.~Milne,
\href{https://doi.org/10.1007/978-1-4613-8655-1_5}{\emph{Abelian varieties}},
in \emph{Arithmetic Geometry} (Storrs, Conn., 1984),
G.~Cornell and J.~H.~Silverman (eds.),
Springer-Verlag, New York, 1986, pp.~103--150.

\bibitem{Milne1986}
J.~S.~Milne,
\href{https://doi.org/10.1007/978-1-4613-8655-1_7}{\emph{Jacobian varieties}},
in \emph{Arithmetic Geometry} (Storrs, Conn., 1984),
G.~Cornell and J.~H.~Silverman (eds.),
Springer-Verlag, New York, 1986, pp.~167--212.

\bibitem{Mochizuki1996}
S.~Mochizuki, \emph{The profinite Grothendieck conjecture for closed
hyperbolic curves over number fields},
J. Math. Sci. Univ. Tokyo \textbf{3} (1996), no.~3, 571--627.

\bibitem{Mochizuki1999}
S.~Mochizuki,
\href{https://doi.org/10.1007/s002220050381}{\emph{The local pro-$p$
anabelian geometry of curves}},
Invent. Math. \textbf{138} (1999), no.~2, 319--423.

\bibitem{Mochizuki2003}
S.~Mochizuki,
\emph{Topics surrounding the anabelian geometry of hyperbolic curves},
in \emph{Galois Groups and Fundamental Groups},
L.~Schneps (ed.), Mathematical Sciences Research Institute Publications,
vol.~41, Cambridge University Press, Cambridge, 2003, pp.~119--165.

\bibitem{Nakamura1990}
H.~Nakamura,
\href{https://doi.org/10.1515/crll.1990.411.205}{\emph{Galois rigidity
of the \'etale fundamental groups of punctured projective lines}},
J. Reine Angew. Math. \textbf{411} (1990), 205--216.

\bibitem{Raynaud1970}
M.~Raynaud,
\href{https://doi.org/10.1007/BF02684651}{\emph{Sp\'ecialisation du
foncteur de Picard}},
Inst. Hautes \'Etudes Sci. Publ. Math. \textbf{38} (1970), 27--76.

\bibitem{Stacks}
The Stacks Project Authors, \emph{The Stacks Project},
\url{https://stacks.math.columbia.edu}, 2026.

\bibitem{Tamagawa1997}
A.~Tamagawa,
\href{https://doi.org/10.1023/A:1000114400142}{\emph{The Grothendieck
conjecture for affine curves}},
Compositio Math. \textbf{109} (1997), no.~2, 135--194.

\bibitem{Tate1967}
J.~T.~Tate,
\href{https://doi.org/10.1007/978-3-642-87942-5_12}{\emph{$p$-Divisible groups}},
in \emph{Proceedings of a Conference on Local Fields}
(Driebergen, 1966), T.~A.~Springer (ed.),
Springer-Verlag, Berlin, 1967, pp.~158--183.

\bibitem{Tsujimura2023}
S.~Tsujimura,
\href{https://doi.org/10.1016/j.aim.2023.109380}{\emph{Birational anabelian
Grothendieck conjecture for curves over arbitrary cyclotomic extension
fields of number fields}},
Adv. Math. \textbf{435} (2023), Paper No.~109380.

\bibitem{Tsujimura2026}
S.~Tsujimura, \emph{Generalizations of birational anabelian Grothendieck
conjecture for curves~I},
RIMS Preprint No.~1997, Research Institute for Mathematical Sciences,
Kyoto University, July 2026.

\end{thebibliography}
\end{document}
